% !TeX root = ../notes.tex
% Proposed prose version.
% This is an editable alternative chapter draft aligned to 05-model-plan.tex.

\chapter{Modelling chemical reactivity}

\begin{notebox}
  \keypoint{What do you need to learn from this chapter?}
  Molecular descriptors become useful for prediction only when they are connected to a defined chemical response, a defensible dataset, and an appropriate model.
  This chapter develops the physical-organic idea of relating structure to reactivity or selectivity, then asks how far a fitted relationship can be interpreted and trusted.
\end{notebox}

\section{From describing molecules to modelling chemistry}

Chapter~4 represented molecules as points in chemical space.
That description did not require a measured chemical outcome because the objective was to compare molecules, visualise their relationships, and choose representative examples.
A reactivity model introduces a response such as rate, conversion, yield, regioselectivity, or enantioselectivity.

For $n$ experiments described by $p$ features, the inputs can be collected in a descriptor matrix $\mathbf{X}$ and the measured responses in a vector $\mathbf{y}$.
A fitted model seeks a relationship
\begin{equation}
  \hat{\mathbf{y}}=f(\mathbf{X}).
\end{equation}
The hat on $\hat{\mathbf{y}}$ matters because the model produces predictions rather than the experimental observations themselves.

\begin{conceptbox}
  \keypoint{Representation, response, and model are separate choices.}
  Changing any one of them can change the apparent success of the modelling exercise and the chemical conclusions drawn from it.
\end{conceptbox}

\section{Physical-organic relationships as models}

Physical-organic chemistry has long used numerical relationships to connect molecular structure with chemical behaviour.
The Hammett relationship is a simple example because one substituent parameter, $\sigma_{\mathrm{X}}$, is related to a logarithmic rate or equilibrium response through the reaction constant $\rho$.
The model is useful not because every substituted reaction lies perfectly on one line, but because the sign, magnitude, and failures of the relationship can be interpreted chemically.

Modern quantitative structure--selectivity relationships extend the same idea to several descriptors,
\begin{equation}
  y=\beta_0+\beta_1x_1+\beta_2x_2+\cdots+\beta_px_p+\varepsilon.
\end{equation}
The fitted coefficients describe how the response changes within the sampled dataset when the associated descriptors vary.
They do not automatically establish that those descriptors are causal variables in the reaction mechanism.

\section{Data, representation, and algorithm}

A modelling project should begin by defining the desired domain rather than by selecting an algorithm.
The domain specifies which substrates, catalysts, reagents, conditions, and response values the model is expected to cover.
The experimental budget, measurement method, and expected precision then constrain which parts of that domain can be sampled.
Raghavan and co-workers organise these decisions into a dataset-design workflow that moves from task definition and experimental constraints to deliberate component selection, data collection, and modelling.\mcite{Raghavan2023}

Chemical datasets are not interchangeable.
High-throughput experiments may provide many internally consistent measurements over a deliberately chosen region, but they can remain chemically narrow.
Published reaction scopes contain chemically curated examples, but often contain few unsuccessful experiments and may change several structural features at once.
Literature databases increase scale at the cost of heterogeneous procedures, analytical methods, and reporting conventions.

The response distribution also affects the sensible modelling task.
A dataset containing a large population of inactive reactions and a smaller continuum of productive reactions may be treated as a classification problem, a regression problem, or a staged combination of both.
Haas, Kalyani, and Sigman show how response distributions, sparse data, negative results, and validation choices should inform the modelling strategy.\mcite{Haas2025}

\section{Interpretable multivariate models}

Multiple linear regression is attractive because its mathematical form is transparent and its descriptors can be chemically motivated.
Interpretability nevertheless depends on more than being able to read a coefficient table.
If two descriptors carry nearly the same information, their fitted coefficients may be unstable even when the combined predictions remain satisfactory.
Scaling, descriptor selection, repeated fitting, and inspection of correlations are therefore part of chemical interpretation rather than optional preprocessing.

A fitted model must be tested on data that were not used to determine its parameters.
A random row split is useful only when the intended future prediction resembles another random row from the same population.
If the scientific question concerns a new substrate family, catalyst scaffold, or reagent class, then the corresponding component should be held out during validation.

Residuals are the differences between observed and predicted responses.
Small, unstructured residuals are consistent with a model capturing the dominant relationship over the sampled domain.
A group of residuals associated with one substrate class, catalyst family, or condition range can instead indicate a missing descriptor, a measurement problem, or a change in chemical regime.

\begin{insightbox}
  \keypoint{A residual can be chemically informative.}
  The question is not only how large the error is, but whether the error has chemical structure.
\end{insightbox}

\section{Reactivity cliffs and mechanistic regimes}

A smooth model assumes that similar inputs tend to produce similar responses.
Reactivity cliffs violate that expectation because a small change in representation is associated with a large change in chemical outcome.
Such cliffs can expose a boundary between catalytic states or mechanisms that is compressed into one descriptor space.

Newman-Stonebraker and co-workers combined data from nickel- and palladium-catalysed cross-coupling reactions with descriptors from the \textit{kraken} phosphine library.\mcite{NewmanStonebraker2021}
Minimum buried volume separated active and inactive ligand regimes close to a threshold, whereas cone angle did not provide the same useful classification.
Spectroscopic studies then connected the descriptor boundary to phosphine ligation state at nickel.
The case is powerful because the statistical separation is supported by an experimentally testable chemical interpretation.

\begin{warningbox}
  A sharp statistical boundary should prompt a chemical question.
  It should not be labelled as a mechanistic boundary until independent evidence supports that interpretation.
\end{warningbox}

\section{Predicting catalyst selectivity}

\subsection{Designed catalyst space: Denmark}

Zahrt and co-workers began from a large in-silico library of synthetically accessible chiral phosphoric-acid catalysts.\mcite{Zahrt2019}
They calculated conformer-aware average steric occupancy descriptors and selected a representative training subset from the resulting catalyst space.
The training set was therefore designed to cover the representation rather than assembled only from familiar catalysts.

The study then tested predictions for new substrate combinations and catalysts outside the experimentally observed training combinations.
This distinction is important because a predicted response can lie beyond the selectivity values observed during training while the corresponding catalyst remains inside the sampled descriptor domain.
Response-range extrapolation and chemical-space extrapolation are related but not identical risks.

\subsection{Nonlinear quantum-descriptor model: Sunoj}

Singh and co-workers assembled asymmetric-hydrogenation data across several catalyst families and represented catalysts and substrates using quantum-derived descriptors.\mcite{Singh2020}
Random-forest models captured nonlinear relationships that would be difficult to express through a small linear equation.
Performance decreased as the external test set became more chemically distinct from the training catalyst families, illustrating that chemical novelty is a graded rather than binary concept.

The authors also used a decision tree to identify descriptors associated with high predicted selectivity.
Those branches provide hypotheses about useful combinations of catalyst and substrate properties.
They do not by themselves prove which interactions control the transition state.

\section{When should a prediction be trusted?}

A model's applicability domain is the region over which there is evidence that its predictions remain reliable.
The domain depends on the sampled chemistry, the representation, the response, and the model.
It cannot be established from an algorithm name alone.

Interpolation estimates responses between well-supported regions of the training data.
Extrapolation places the prediction in a region where the model has little or no direct support.
A point can appear numerically close to training data while belonging to a different mechanism, and a point can appear distant because the chosen descriptors emphasise irrelevant structural differences.
Applicability is therefore a chemical judgement supported by numerical diagnostics rather than a single universal distance threshold.

Validation should reproduce the novelty expected in use.
Random reaction splits estimate performance for unseen rows drawn from the same mixture of components.
Component-held-out, scaffold-held-out, or time-split validation can test stronger claims about unfamiliar chemistry.
The workshop returns to this distinction through the rule that an unseen experiment is not necessarily unseen chemistry.

\section{Prediction is not explanation}

Predictive success shows that a relationship is useful for estimating a response over a tested domain.
Mechanistic explanation requires evidence that the proposed variables and causal story survive intervention and independent measurement.
A descriptor may be important because it correlates with the true controlling variable, because it identifies a family of molecules, or because it leaks information about the response.

\begin{warningbox}
  \keypoint{Feature importance is not automatically mechanistic importance.}
  A model can return a precise numerical prediction even when the chemistry provides no reason to trust it.
\end{warningbox}

The most productive modelling studies use prediction and chemical interpretation together.
They report the domain, validation design, residuals, and failures, and they use surprising predictions to motivate new experiments.
This prepares the transition to Chapter~6, where the central problem is no longer only what the model predicts, but which experiment should be performed next.

\section{Summary}

Molecular representations become model inputs when they are paired with a defined chemical response.
Physical-organic relationships provide a chemistry-led route from one-variable correlations to interpretable multivariate models.
Dataset design, response distribution, validation, and applicability domain determine which prediction claims are defensible.
Residuals and reactivity cliffs can reveal missing chemistry or regime changes.
Predictive performance supports use within a tested domain, but it does not by itself establish a mechanism.
